\cotefile{82}
\scanpage{2}
\margin{(29 p)}

Algèbres de Lie de rang 3 associées aux formes bil. symétriques de rang 3

% archive sheet number: 1
\scanpage{3}
\fbox{Formes quadratiques de rang trois et \unsure{algèbre|algèbres} de Lie}

$S$ schéma, $E,E'$ Modules vectoriels de rang 3 sur $\underline S$, en dualité par accouplement
\[
 E\times E'\longrightarrow\underline{\mathrm O}_S\qquad (x,x')\longmapsto\langle x,x'\rangle \tag*{(1)}
\]
d'où accouplement
\[
 \underline{\det}(E)\otimes\underline{\det}(E')\simeq\underline{\mathrm O}_S. \tag*{(2)}
\]
(où on pose
\[
 \underline{\det} E\simeq\bigwedge^3 E\qquad )
\]

Notons que les accouplements
\[
\begin{aligned}
 E\times\bigwedge^2 E&\longrightarrow\bigwedge^3 E=\underline{\det} E
 &&(\deleted{\unsure{x\mapsto(y\wedge z)}})\quad
 \inserted{(x,\omega)\mapsto x\wedge\varphi=\unsure{\varphi\wedge x}}\\
 E'\times\bigwedge^2 E'&\longrightarrow\bigwedge^3 E'=\underline{\det} E'
 &&(x',\varphi')\mapsto x'\wedge\varphi'\inserted{=\varphi'\wedge x'}
\end{aligned}
\]
sont aussi dualisants, donc on a des isomorphismes
\[
\left\{\begin{aligned}
 \bigwedge^2 E&\xrightarrow{\overset{\alpha_0}{\sim}}\check E\otimes\underline{\det} E\simeq E'\otimes\underline{\det}(E)\\
 \bigwedge^2 E'&\xrightarrow{\overset{\alpha'_0}{\sim}}\check E'\otimes\underline{\det} E'\simeq E\otimes\underline{\det}(E')
\end{aligned}\right. \tag*{(3)}
\]

Supposons maintenant $E$ ``spécial'' i.e. muni d'une \uline{base}
\[
 \omega\in\Gamma(S,\underline{\det}(E)) \tag*{(4)}
\]
définissant donc un isomorphisme
\[
 \underline{\mathrm O}_S\xrightarrow{\sim}\det(E)\qquad
 \deleted{\unsure{\lambda}}\ \lambda\mapsto\lambda\omega \tag*{(5)}
\]
et une base duale
\[
 \omega'=\omega^{-1}\in\Gamma(S,\underline{\det}(E')) \tag*{(\unsure{4|5}$'$)}
\]
définissant
\[
 \underline{\mathrm O}_S\xrightarrow{\sim}\det(E')\qquad\lambda\mapsto\lambda\omega' \tag*{(5$'$)}
\]
Compte tenu des iso. (5) (5$'$), les isom. (3) définissent des iso.
\[
\left\{\begin{aligned}
 \bigwedge^2 E&\xrightarrow{\overset{\alpha}{\sim}}E'\\
 \bigwedge^2 E'&\xrightarrow{\overset{\alpha'}{\sim}}E\deleted{\unsure{'}}
\end{aligned}\right.\qquad\text{définis donc par} \tag*{(6)}
\]
\[
 \langle x,\alpha(\varphi)\rangle\omega=x\wedge\varphi,\quad
 \deleted{\unsure{\langle x',\alpha'(\varphi')\rangle\omega'=}}\;
 \langle\alpha'(\varphi'),x'\rangle\unsure{\omega|\omega'}=\varphi'\wedge x' \tag*{(7)}
\]

% archive sheet number: 2
\scanpage{4}
Si $E$ est muni d'une base
\[
 (e_1,e_2,e_3)\ \unsure{\in|e}\qquad (e_i\in\Gamma(S,E)) \tag*{(8)}
\]
définissant la base duale
\[
 (e'_1,e'_2,e'_3)\qquad(e'_i\in\Gamma(S,E')) \tag*{(8$'$)}
\]
par
\[
 \langle e_i,e'_j\rangle=\delta_{ij}\qquad\text{(symb. de Kronecker)}
\]
on a des bases de $E,E'$ par
\[
\left\{\begin{aligned}
 \bar e_1=e_2\wedge e_3,\ \bar e_2=e_3\wedge e_1,\ \bar e_3=e_1\wedge e_2
 &\qquad(\text{base de }\bigwedge^2 E)\\
 \bar e'_1=e'_2\wedge e'_3,\ \bar e'_2=e'_3\wedge e'_1,\ \bar e'_3=e'_1\wedge e'_2
 &\qquad(\text{--- }\bigwedge^2 E')
\end{aligned}\right. \tag*{(\unsure{9|8})}
\]
et on a
\[
 \left\{\begin{aligned}
 \alpha(\bar e_i)&=e'_i\\
 \alpha'(\bar e'_i)&=e_i
 \end{aligned}\right. . \tag*{(\unsure{10|9})}
\]
Notons \deleted{\unsure{donc}}

\uline{Observation} Si $E$ est un Module vectoriel ``spécial'' de rang 3, alors on a \deleted{une} bijection\deleted{\unsure{s}} canoniques
\diagram{\begin{tikzcd}[ampersand replacement=\&,column sep=small]
{\substack{\text{Formes bilinéaires}\\ E\times E\to\underline{\mathrm O}_S}}
 \arrow[rr,"\sim"] \arrow[dr,"\sim"] \& \&
 {\substack{\text{structures d'Algèbre}\\ \text{alternée sur } E'}}\\
 \unsure{\text{---}} \& \operatorname{Hom}_{\underline{\mathrm O}_S}(E,E')\arrow[ur,"\sim"] \&
\end{tikzcd}}
où $\overset{\sim}{\searrow}$ \uline{associe} à $\underline{\varphi(x,y)}$ \uline{l'application $f$}$:E\to E'$ définie par
\[
 f(x)(y)\bigl(=\langle y,f(x)\rangle\bigr)=\varphi(x,y) \tag*{(11)}
\]
et où $\overset{\sim}{\nearrow}$ est l'application qui à $\deleted{\unsure{\varphi}}f$ associe la structure d'algèbre alternée
\[
 (x',y')\longmapsto [x',y']_\varphi=\deleted{\unsure{\varphi}}f\alpha'(x'\wedge y')\ . \tag*{(12)}
\]

% archive sheet number: 3
\scanpage{5}
que l'effet d'une permutation circulaire des vecteurs de base sur la matrice $\varphi$ s'obtient par une permutation circulaire sur les \inserted{trois groupes de} lettres $(a,b,c)$, $(p,q,r)$, $(p',q',r')$. Ceci posé, les \deleted{\unsure{la}} formules (13) (13 bis) se récrivent
\[
\left\{\begin{aligned}
 [[e'_1\wedge e'_2],e'_3]&=(ap'-qr')e'_1+(p'r-bq)e'_2+(p'q'-pq)e'_3\\
 [[e'_2\wedge e'_3],e'_1]&=(q'r'-qr)e'_1+(bq'-rp')e'_2+(q'p-cr)e'_3\\
 [[e'_3,e'_1],e'_2]&=(r'q-ap)e'_1+(r'p'-rp)e'_2+(cr'-pq')e'_3
\end{aligned}\right. \tag*{(16)}
\]
d'où
\[
\begin{aligned}
 \text{Jác}(e_1,e_2,e_{\unsure{1|3}})&=[[e_1,e_2],e_3]+[[e_2,e_3],e_1]+[[e_3,e_1],e_2]\\
 &=[a(p'-p)+(q'r'-qr)]e'_1+[b(q'-q)+(r'p'-rp)]e'_2+[c(r'-r)+(p'q'-pq)]e'_3
\end{aligned}\tag*{[17]}
\]
Donc, pour que l'Algèbre $E'_\varphi$ ($E$ muni de $[\, ,\,]_\varphi$) soit une Algèbre de Lie, il faut et il suffit qu'elle satisfasse les conditions
\[
\left\{\begin{aligned}
 a(p'-p)+(q'r'-qr)&=0 &\text{ou}\quad (ap'-qr)-(ap-q'r')&=0\\
 b(q'-q)+(r'p'-rp)&=0 &\text{ou}\quad (bq'-rp)-(bq-r'p')&=0\\
 c(r'-r)+(p'q'-pq)&=0 &\text{---}\quad(cr'-pq)-(cr-p'q')&=0
\end{aligned}\right. \tag*{(18)}
\]

\uline{NB} Avec la notation évidente pour les mineurs d'ordre 2 $M_{\substack{\unsure{i|1|k}\\\unsure{j|3|\ell}}}^{\unsure{k\ell|12|ij}}(\varphi)$ d'une matrice $\varphi$, les premiers membres des équations (18), i.e. les coefficients de l'alternateur jacobien (18), sont
\[
\left\{\begin{aligned}
 &\operatorname{Min}_{\substack{1\\3}}^{12}(\varphi)-\operatorname{Min}_{\substack{1\\3}}^{12}({}^{t}\varphi)\\
 &\operatorname{Min}_{\substack{2\\1}}^{23}(\varphi)-\operatorname{Min}_{\substack{2\\1}}^{23}({}^{t}\varphi)\\
 &\operatorname{Min}_{\substack{3\\2}}^{31}(\varphi)-\operatorname{Min}_{\substack{3\\2}}^{\unsure{31|21}}({}^{t}\varphi)
\end{aligned}\right. \tag*{(19)}
\]

% archive sheet number: 4
\scanpage{6}
En termes des bases (8) (8$'$), $\varphi$ sera donnée par une matrice (notée \deleted{\unsure{aussi} $\varphi$ \unsure{$\underline{\mathrm O}_S$}})
\[
 \varphi=(\varphi_{ij})_{1\leq i,j\leq3}\qquad \varphi_{ij}=\varphi(e_i,e_j) \tag*{(11)}
\]
-- la matrice de $f$ est alors \uline{transposée} de la matrice $\varphi$ -- et on trouve que la structure d'algèbre alternée correspondante est donnée par
\[
 [e'_2\unsure{\wedge|,}e'_3]_\varphi=\varphi\alpha'(e'_2\wedge e'_3)=\varphi\alpha'(\bar e'_1)
 \underset{(10)}{=}\varphi(e_1)=\sum_j\varphi_{1j}e'_j,
\]
et de \deleted{$[e'_3\wedge e'_1]=$} même pour les deux autres \deleted{produits}\inserted{crochets}
\[
\left\{\begin{aligned}
 [e'_2,e'_3]_\varphi&=\sum_j\varphi_{1j}e'_j\\
 [e'_3,e'_1]_\varphi&=\sum_j\varphi_{2j}e'_j\\
 [e'_1,e'_2]_\varphi&=\sum_j\varphi_{3j}e'_j
\end{aligned}\right. \tag*{(12)}
\]
Ceci posé, on trouve (laissant tomber l'indice $\varphi$ dans les \inserted{crochets})
\[
 [[e'_1\wedge e'_2],e'_3]=\sum_j(\varphi_{32}\varphi_{1j}-\varphi_{31}\varphi_{2j})e'_j \tag*{(13)}
\]
et de même pour les expressions \unsure{déduite|déduites} par permutations circulaires
\[
\left\{\begin{aligned}
 [[e'_2\wedge e'_3],e'_1]&=\sum_j(\varphi_{13}\varphi_{2j}-\varphi_{12}\varphi_{3j})e'_j\\
 [[e'_3\wedge e'_1],e'_2]&=\sum_j(\varphi_{21}\varphi_{3j}-\varphi_{23}\varphi_{1j})e'_j
\end{aligned}\right. \tag*{(13 bis)}
\]
Les \deleted{coefficients} symétries des formules apparaissent mieux en posant
\[
 \deleted{\varphi_{11}=a,\ \unsure{\varphi_{1\cdot}|\varphi}=}
\]
\[
\left\{\begin{aligned}
 \varphi_{11}=a,\quad &\varphi_{22}=b,\quad \varphi_{33}=c\\
 \varphi_{23}=p,\quad &\varphi_{31}=q,\quad \varphi_{12}=r\\
 \varphi_{32}=p',\quad &\varphi_{13}=q',\quad \varphi_{21}=r'
\end{aligned}\right. \tag*{(14)}
\]
i.e. en écrivant
\[
 \varphi=\begin{pmatrix}a&r&q'\\ r'&b&p\\ q&p'&c\end{pmatrix} \tag*{(15)}
\]
de sorte que la matrice transposée de $\varphi$ s'obtient en échangeant $p$ avec $p'$, $q$ avec $q'$, $r$ avec $r'$, et

% archive sheet number: 5
\scanpage{7}
On trouve en particulier

\uline{Proposition} Les formes bilinéaires \uline{symétriques} $\varphi$ sur $E$ définissent des structures d'\uline{Algèbre de Lie} sur $E'$.

\uline{Exemple} Prenons $E$ muni de la base (8), \deleted{\unsure{et envi-}}\inserted{d'où} identification
\[
 E\simeq\underline{\mathrm O}_S^3 \tag*{(20)}
\]
\unsure{écrivons|décrivons|choisissons|dirons} sur $E$ la forme quadratique
\[
 q(x,y,z)=xy+z^2 \tag*{(21)}
\]
correspondant à la forme bilinéaire associée
\[
 \varphi((x,y,z)(x',y',z'))=(xy'+yx')+2zz' \tag*{(22)}
\]
donc
\[
\left\{\begin{aligned}
 q(e_1)=q(e_2)&=0,\quad q(e_3)=1\\
 \varphi(e_1,e_2)&=1,\quad \varphi(e_2,e_3)=\varphi(e_3,e_1)=0\\
 \varphi(e_1,e_1)&=\varphi(e_2,e_2)=0\qquad\varphi(e_3,e_3)=2
\end{aligned}\right. \tag*{(23)}
\]
\margin{\uline{NB} $\varphi(\xi,\xi)=2q(\xi)$}
(ou encore
\[
\left\{\begin{aligned}
 a=b&=0,\quad c=1\\
 p=p'=q=q'&=0\qquad r=r'=1
\end{aligned}\right.,\quad \text{matrice})\ ,\varphi=\begin{pmatrix}0&1&0\\1&0&0\\0&0&2\end{pmatrix} \tag*{(23 bis)}
\]
Les formules (12) deviennent alors \deleted{\unsure{ainsi}}
\[
\left\{\begin{aligned}
 [e'_2\unsure{\wedge|,}e'_3]&=e'_2\\
 [e_3,e'_1]&=e'_1\\
 [e'_1,e'_2]&=2e'_3
\end{aligned}\right. \tag*{(24)}
\]
\margin{NB $\begin{aligned}f(e_1)&=e'_2\\ f(e_2)&=e'_1\\ f(e_3)&=2e'_3\end{aligned}$}
ou encore, posant
\[
 X=e'_1,\quad Y=e'_2,\quad H=e'_3 \tag*{(25)}
\]
on reconnait les formules familières
\[
 [H,X]=X,\quad[H,Y]=-Y,\qquad[X,Y]=2H \tag*{(26)}
\]
de l'algèbre de Lie \deleted{\unsure{du|de}} ``universelle'' (sur les entiers) de $\mathrm{GP}(1)$.

% archive sheet number: 6
\scanpage{8}
Notons que pour une forme bilinéaire
\[
 \varphi(x,y)=\sum_{1\leq i,j\leq3}\varphi_{ij}x_i y_j
\]
on a le discriminant de $\varphi$ (ou de la forme quadratique $q$, si $\varphi$ est associé à une telle forme)
\[
 \delta(\varphi)=\unsure{-|}\det(\varphi_{ij})\,\omega'^{\unsure{\otimes}2}
 \qquad(\text{où }\omega'=e'_1\wedge e'_{\unsure{1|2}}\wedge e'_3), \tag*{(27)}
\]
et si $\varphi$ est associé à une forme quadratique
\[
 \inserted{q(x,y,z)=}ax^2+by^2+cz^2+pyz+qzx+rxy \tag*{(28)}
\]
de sorte que la matrice de $\varphi$ est
\[
 \varphi=\begin{pmatrix}2a&r&q\\r&2b&p\\q&p&2c\end{pmatrix} \tag*{(29)}
\]
on a
\[
 \delta(q)=\delta(\varphi)=2(4abc+pqr-ap^2-bq^2-cr^2)\,\omega'^{\otimes2} \tag*{(30)}
\]
et le ``discriminant divisé'' $\delta'(q)$ de $q$ est donné par
\[
 \delta'(q)=(4abc+pqr-ap^2-bq^2-cr^2)\,\omega'^{\otimes2}, \tag*{(31)}
\]
\deleted{ce qui} Dans le cas de la forme $q$ de (21) d'où (23 bis), on trouve donc
\[
 \delta'(q)=-\omega'^{\otimes2}\qquad(\omega'=e'_1\wedge e'_2\wedge e'_3) \tag*{(32)}
\]

En s'inspirant de cette circonstance, on définit une structure \uline{spéciale orthogonale} sur un Module \deleted{\unsure{quadrati-|quadratique}} vectoriel de rang 3, comme la donnée d'un couple $(q,\omega)$, où $q$ est une forme quadratique sur $E$, $\omega$ une base de $\underline{\det} E$, satisfaisant

% archive sheet number: 7
\scanpage{9}
\[
 \delta'(q)=-\omega^{\otimes(-2)} . \tag*{(33)}
\]
\uline{NB} Si $E$ était de rang $2n$ ou $2n+1$, on \deleted{\unsure{satisfa-|satisferait}} exigerait
\[
 \delta'(q)=(-1)^n\omega^{\otimes(-\unsure{2})} \tag*{(34)}
\]
de sorte que pour la forme \unsure{standard|standart} sur $\mathbb Z$
\[
 q(x)=\sum_1^{\unsure{n}}x_{2i}\unsure{y|x}_{2i+1}
 \qquad\text{resp.}\quad q(x)=\sum_1^n x_{2i}x_{2i+1}+x_{2n+1}^2 \tag*{(35)}
\]
dont le discriminant resp. le discriminant divisé est
\[
 \delta'(q)=(-1)^n\omega^{\otimes(-2)}\qquad \omega=e_1\wedge\text{---}\wedge e_N \tag*{(36)}
\]
($N=2n$ ou $2n+1$), on puisse trouver effectivement une structure spéciale orthogonale associée. On note que, dans le cas général, il existe \inserted{loc. fppf} une structure sp. orthogonale associée à $q$ ssi $q$ est lisse (i.e. $\delta'(q)$ est une base de $\det.(E)^{\otimes(-2)}$), et que les structures en question, pour base variable, forment un \uline{torseur} sous $\unsure{\mu_2|\boldsymbol{\mu}_2}$.

\begin{center}\rule{1.3cm}{0.4pt}\end{center}

Les développements précédents sur les structures d'Algèbre alternées sur $E'$ associées aux formes bilinéaires sur $E$ (i.e. aux applications linéaires
\[
 f:E\longrightarrow E')
\]
peuvent s'appliquer en échangeant les \unsure{rôles|rôle} de $E,E'$, de sorte qu'on associe une structure d'Algèbre alternée \inserted{sur $E$} à tte forme bil. $\psi$ sur $E'\times E'$, i.e. à toute appl. lin.
\[
 f':E'\longrightarrow E .
\]

% archive sheet number: 8
\scanpage{10}
Or, si on part de $\varphi$ i.e de
\[
 f:E\longrightarrow E' \tag*{(36)}
\]
on peut lui associer
\[
 \bigwedge^2 f:\bigwedge^2 E\longrightarrow\bigwedge^2 E' \tag*{(37)}
\]
qui en vertu des isom. (6) définit une application linéaire
\[
 f'=\alpha'\circ\bigwedge^2 f\circ\alpha^{-1}:E'\longrightarrow E \tag*{(38)}
\]
\deleted{\unsure{Du point|De sorte}} Par définition, on aura
\[
 f'(\alpha(\deleted{\varphi}\unsure{x'|x}\wedge\unsure{y'|y}))=\alpha'(f(x)\wedge f(y))\qquad x,y\in\Gamma(S,E) \tag*{(39)}
\]
i.e.
\[
 \langle f'(\alpha(x\wedge y)),x'\rangle\omega'=f(x)\wedge f(y)\wedge x'\qquad(x'\in\Gamma(S,E')\unsure{)} \tag*{(40)}
\]
En fait, en termes d'une base \inserted{(8) et de sa duale}, les coefficients de la matrice \inserted{exprimant} $f'$ sont les 9 mineurs de la matrice de $f$ -- ce qui montre en particulier que si la forme \inserted{bil.} $\varphi$ associée à $f$ est symétrique, il en est de $=$ de la forme $\varphi'$ associée à $f'$. Dans ce cas, la structure d'algèbre de $E$ associée \deleted{$[\unsure{x,y|\, ,\,}]_{\varphi'}$} (qui n'est autre que
\[
 [x,y]_{\varphi'}=\alpha'\bigwedge^2 f(x\wedge y)=\alpha'(f(x)\wedge f(y)) \tag*{(40 bis)}
\]
\unsure{explicité|explicit.} par (40), qui s'écrit
\[
 \langle[x,y]_{\varphi'},x'\rangle\omega'=f(x)\wedge f(y)\wedge x'\qquad ) \tag*{(40 \deleted{bis}\inserted{ter})}
\]
est \inserted{encore} une structure d'\uline{algèbre de Lie}.

% archive sheet number: 9
\scanpage{11}
\uline{Remarques} Des \deleted{formules de}\inserted{on pose}
\[
 \underset{\text{(discriminant)}}{\delta(\varphi)}=\Delta\omega^{\unsure{\otimes}{-2}},\qquad\Delta\in\Gamma(S,\underline{\mathrm O}_S) \tag*{(41)}
\]
et les formules de Kramer donnent alors
\[
 \left\{\begin{aligned}
 f'f&=\Delta\,\mathrm{id}_E\\
 ff'&=\Delta\,\mathrm{id}_{E'}
 \end{aligned}\right. \tag*{(42)}
\]
On se gardera de croire que les relations entre $f,f'$ sont symétriques. P.ex. si $f\unsure{^{a}|}$ est de rang $\leq1$ (i.e. les mineurs de rang 2 de $f$ sont nuls) $f'$ est \uline{nul}, et réciproquement, tandis que si $f$ est nul, $f'$ est nul lui-$=$ (et non seulement son rang est $\leq1$)!

Notons que si $\omega$ est multiplié par $\lambda\in\Gamma(S,\underline{\mathrm O}_S^{\unsure{*|\times}})$, ($\varphi$ donc $f$ restant invariant) $\Delta$ et $f'$ sont multipliés l'un et l'autre par $\lambda^2$ -- tandis que le crochet $[\, ,\,]_\varphi$ \inserted{sur $E'$} est multiplié par $\lambda$, et \inserted{le crochet} $[\, ,\,]_{\varphi'}$ sur $E'$, pour \deleted{\unsure{et}} une $\varphi'$ fixée (indépendamment d'une $\varphi$ sur $E$) est multiplié (comme $\omega'=\omega^{\otimes(-1)}$) par $\lambda^{-1}$.

Mais si $\varphi'$ est la forme associée à $f$ via $f'$, comme \inserted{$f'$ donc} $\varphi'$ est \deleted{multiplié}\inserted{remplacé} par $\lambda^2\varphi'$, $[\, ,\,]_{\varphi'}^{\unsure{\omega|}}$ est remplacé par $\lambda[\, ,\,]_{\varphi'}$ : Ainsi $[\, ,\,]_\varphi$ \inserted{sur $E$} et $[\, ,\,]_{\varphi'}$ sur $E'$, associés à un couple

% archive sheet number: 10
\scanpage{12}
$(\varphi,\omega)$ ($\varphi$ forme bilinéaire sur $E$, $\omega$ base de $\underline{\det} E$) se trouvent multipliés par un $=$ facteur $\lambda$, quand $\omega$ est multiplié par $\lambda$\inserted{ et $\varphi$ \unsure{reste|restant} \unsure{fixée|fixe}}. D'autre part, pour $\omega$ fixé et $\varphi$ variable, le crochet $[\, ,\,]_{\varphi\deleted{\unsure{'}}}$ sur $E'$ dépend linéairement de $\varphi$, tandis que le crochet $[\, ,\,]_{\varphi'}$ sur $E$ (tout comme $f'$ lui-$=$ dont il est déduit linéairement) dépend \uline{quadratiquement} de $\varphi$ ; ainsi, si $\varphi$ est \deleted{multiplié}\inserted{remplacé} par $\mu\varphi$ ($\mu\in\Gamma(S,\underline{\mathrm O}_S)$) $[\, ,\,]_\varphi$ est multiplié par $\mu$, $[\, ,\,]_{\varphi'}$ par $\mu^2$....

Échangeant les rôles de $E,E'$, on voit de $=$ qu'à $f':E'\to E$ est associée une $(f')'=f'':E\to E'$ -- mais en général celle-ci \uline{n'est pas} $f$. De façon précise, si \deleted{$\delta(\varphi')=$}
\[
 \delta(\varphi')=\Delta'\mathbin{.}\omega^{\otimes2}\qquad
 \Delta'\in\Gamma(\unsure{S'|S},\underline{\mathrm O}_{\unsure{S'|S}}) \tag*{(43)}
\]
on aura
\[
 \left\{\begin{aligned}
 f'f''&=\Delta'\,\mathrm{id}_E\\
 f''f'&=\Delta'\,\mathrm{id}_{E'}
 \end{aligned}\right. \tag*{(44)}
\]
ce qui \deleted{\unsure{sugg-}} implique, si $\Delta$ est une unité i.e. $f$ un iso, donc $\bigwedge^2f$ et $f'$ aussi, donc $\Delta'$ une unité [et inversement, car si $f'$ est un iso les formules (42)

% archive sheet number: 11
\scanpage{13}
montrent que $f$ l'est aussi...], alors \unsure{le|la|là} \deleted{\unsure{d'}}on trouve
\[
 f'=\Delta f^{-1},\quad f''=\Delta' f'^{-1}=(\Delta'\Delta^{-1})f \tag*{(45)}
\]
ce qui s'écrit
\[
 \Delta f''=\Delta' f \tag*{(45 bis)}
\]
(d'après le ``principe de conservation des identités algébriques'', cette \deleted{\unsure{$f$}} identité reste vraie sans la condition \deleted{\unsure{sur}} $\Delta$ une unité).

De plus, dans le cas de validité de (45) ($\Delta$ unité) on en tire
\[
 \det(f')=\Delta^3\det(f)^{-1}=\Delta^3(\Delta\omega^{\unsure{\otimes}{-2}})^{-1}
\]
i.e.
\[
 \det f'=\Delta^2\omega^{\otimes2} \tag*{(\unsure{47|46})}
\]
\[
 \Delta'=\Delta^2\qquad\text{i.e.} \tag*{\deleted{(\unsure{48|47})}}
\]
\deleted{\unsure{qui reste}} de sorte que (45 bis) devient
\[
 f''=\Delta f \tag*{(48)}
\]
D'après le $=$ principe de conservation des identités Algébriques, cette identité reste valable pour \unsure{toute|tte|tout} $f$ (non seulement si $f$ est un isomorphisme). Ceci \deleted{\illegible{1}} \deleted{montre} .... que si $\Delta=1$ i.e. $f$ un iso et $\omega=\det f$, alors la relation entre $f$ et $f'$ est symétrique. Mais \unsure{nous|non},

% archive sheet number: 12
\scanpage{14}
nous intéresserons surtout au cas (21), \unsure{où|ou} $\Delta=2$, sans nous interdire de spécialiser en car. 2 où on aura $\Delta=0$...

Revenons toujours à la situation $(E,\omega,\varphi)$, donnant naissance à $f:E\to E'$, $f':E'\to E$, d'où $\varphi':E'\times E'\longrightarrow\underline{\mathrm O}_S$, enfin à $[\, ,\,]_\varphi$ sur $E'$ et $[\, ,\,]_{\varphi'}$ sur $E'$. On va montrer que $f:E\to E'$ est un hom. d'algèbre, et \inserted{(si $\varphi$ sym)} on va faire opérer l'algèbre de Lie $E'$ sur $E$ par dérivations \inserted{d'algèbre} qui invarient $\varphi$, de façon que $E\xrightarrow{f}E'$ commute avec l'action de l'algèbre de Lie $E'$ sur $E$ et $E'$. Enfin, lorsque $\varphi$ \inserted{(ainsi que $\omega$)} est associé à une forme quadratique lisse spéciale, on va trouver une interprétation de la situation d'algèbres de Lie en termes d'un hom de schémas en groupes lisses $G\to G'$ et d'une opération de $G'$ sur $G$. \deleted{\unsure{Pensons|Prouvons|Prenons}}\inserted{Vérifions} d'abord
\[
 f[x,y]_{\varphi'}=
 \underset{\substack{\parallel\ \text{déf }(\unsure{12|11})\\ f\alpha'(\unsure{fx\wedge fy|f(x)\wedge f(y)})}}{[f(x),f(y)]_\varphi}\ ; \tag*{[49]}
\]
\deleted{en prenant le pr. \unsure{scalaires|scalaire} des deux mbs avec un $\unsure{Z|z}\in\Gamma(S,E)$, la formule devient, et posant $ {}^{t}f(z)=z'$, la formule à vérifier devient}

% archive sheet number: 13
\scanpage{15}
En vertu de la déf. (12), le deuxième mb est $f\alpha'(f(x)\wedge f(y))$, de sorte qu'il suffit de prouver
\[
 [x,y]_{\varphi'}=\deleted{\unsure{\alpha'f(f(x)}}\ \alpha'(f(x)\wedge f(y)) \tag*{(50)}
\]
ce qui n'est autre que la définition (40 bis).

\uline{NB} En général, $f':E'\to E$ n'est pas compatible aux structures d'Algèbres, $=$ si $f$ est un isomorphisme i.e. $\Delta$ une unité. Dans ce cas, on a $f'=\Delta f^{-1}$ (45), \deleted{\unsure{donc}} \unsure{or|} $f^{-1}$ est \deleted{\unsure{un}} isom. d'Algèbres comme inverse d'un isomorphisme, i.e. on a
\[
 f^{-1}([x',y'])=[f^{-1}(x'),f^{-1}(y')]
\]
Si $f'=\Delta f^{-1}$ \unsure{était} un hom d'algèbre, la \deleted{\unsure{même}} formule précédente \unsure{resterait} valable en y remplaçant $f'$ par $\Delta f^{-1}$, i.e. on aurait à la fois $a=b$ et $\Delta a=\Delta^{\unsure{\prime|2}} b$ i.e. $a=\Delta b$, donc $\Delta a=a$. \deleted{Si les} On en conclut aisément (comme les $[x',y']$ engendrent $E'$ localement ...) que $\Delta=1$, \deleted{Donc} i.e. $f'=f^{-1}$. Donc dès que $\Delta\ne1$, $\Delta$ inversible, $f'$ n'est \uline{\unsure{p=|pas}} un hom d'Algèbres....

% archive sheet number: 14
\scanpage{16}
On se propose maintenant, à tout $x'\in\Gamma_S(E')$ d'associer
\[
 \rho_E(x'):E\longrightarrow E \tag*{(51)}
\]
\deleted{dérivation} endomorphisme de $E$ qui invarie $\varphi$ (donc $[\, ,\,]_\varphi$) infinitésimalement :
\[
 \varphi(\rho_E(x')x,y)+\varphi(x,\rho_E(x')y)=0 \tag*{(52)}
\]
ce qui impliquera
\[
 \rho_E(x')[x,y]_{\varphi'}=[\rho_E(x')x,y]_{\varphi'}+[x,\rho_E(x')y]_{\varphi'} \tag*{(53)}
\]
de sorte qu'on obtient un hom
\[
 E'\xrightarrow{\rho_E}\underline{\mathrm{Hom}}^{\varphi}(E,E)
 \quad\text{(endom. de $E$ invariant infinitésimalement $\varphi$)} \tag*{(54)}
\]
qui en fait est un hom d'algèbres de Lie
\[
 \rho_E([x',y']_{\unsure{\varphi'|\varphi}})=
 \underbrace{[\rho_E(x'),\rho_E(y')]}_{\text{crochet d'endomorphismes}}. \tag*{(55)}
\]
\deleted{(56)} et que l'on ait
\[
 f(\rho_E(x')x)=\operatorname{ad}_{[\, ,\,]_\varphi}(x')f(x)
 \left(\overset{\text{déf}}{=}
 \underset{\substack{\text{(12)}\ \longrightarrow\ \parallel\\ f\alpha'(x'\wedge f(x))}}{[x',f(x)]_\varphi}\right) \tag*{(56)}
\]
% margin note written sideways
\margin{On a oublié une formule naturelle, en introduisant l'opération adjointe $\widetilde{\rho}_E$ de $E$ sur lui-$=$, et l'opération $\widetilde{\rho}_{E'}$ déduite de $f:E\to E'$ et de l'op. adj. de $E'$ sur $E'$ \(\bigl(\widetilde{\rho}_E(x)\mathbin{.}y=[x,y]_{\varphi'},\ \widetilde{\rho}_{E'}(x)\mathbin{.}x'=\rho_E(f(x))\mathbin{.}x'=[f(x),x']_\varphi.\) C'est la formule (56 bis) $\rho_E(f(\unsure{x'|x}))=\rho_E(x')$ exprimant que \unsure{les|l'} op. $\rho_E$ ``prolongent'' \deleted{\unsure{les}}\inserted{via $f$.} l'op. $\widetilde{\rho}_E$ évidente de $E$ sur lui-$=$}
i.e. $f:E\longrightarrow E'$, en tant que hom d'algèbres de Lie, commute aux \deleted{\unsure{homo-}}\inserted{opérations} \deleted{\unsure{morphismes}} $\rho_E$ et $\operatorname{ad}_{[\, ,\,]_\varphi}$ de l'Algèbre de Lie $E'$ sur l'une et l'autre.

\uline{NB} Le fait que (52) -- sur toute base -- implique (53), est \unsure{standart|standard} par ``transport de structure'', en interprétant (52) comme

% archive sheet number: 15
